\documentclass[12pt,a4paper]{article}
\usepackage[utf8]{inputenc}
\usepackage[T2A]{fontenc}
\usepackage[russian]{babel}
\usepackage{amsmath,amssymb,amsthm}
\usepackage{booktabs,array}
\usepackage{microtype}
\usepackage{geometry}
\geometry{margin=2.3cm}
\setlength{\emergencystretch}{2em}

\newtheorem{proposition}{Предложение}
\newtheorem{nogocriterion}{Критерий провала}

\title{Гипотеза производных секторных мер\\
Один parent trace и разные типы наблюдения}
\author{}
\date{4 августа 2026 г.}

\begin{document}
\maketitle

\section{Гипотеза}

Вторая радикальная идея состоит в том, что электромагнитный, лептонный, нейтринный и family-сектора не обязаны использовать одну и ту же численную нормировку. Но это допустимо только если разные меры являются производными одного parent object, а не независимыми коэффициентами.

\section{Примитивный multi-trace no-go}

Пусть внутренняя algebra полупроста:
\[
\mathcal A=\bigoplus_{i=1}^k M_{n_i}(\mathbb C).
\]
Всякий положительный trace имеет вид
\[
\tau(a)=\sum_{i=1}^k w_i\operatorname{Tr}_i(a_i),
\qquad
w_i>0.
\]
Даже после условия \(\tau(1)=1\) остаются \(k-1\) свободных относительных весов.

\begin{nogocriterion}[Независимые секторные traces]
Если разные меры вводятся как примитивные traces на независимых central blocks, они являются скрытыми непрерывными параметрами. Такая версия второй гипотезы закрывается отрицательно.
\end{nogocriterion}

Для трёх blocks, условно соответствующих color, weak и singlet-направлениям, остаются две свободные ручки. Это не лучше прежней подгонки.

\section{Один trace, разные reductions}

Допустимая версия использует один parent trace и фиксированные projectors. Для projector \(P\) существуют разные, но не взаимозаменяемые типы наблюдения:
\begin{align}
\tau_{loop}(O)&=\operatorname{Tr}(POP),\\
\tau_{avg}(O)&=\frac{\operatorname{Tr}(POP)}{\operatorname{Tr}P},\\
\tau_{state}(O)&=\langle\psi,O\psi\rangle,
\qquad \|\psi\|=1.
\end{align}
Первый trace суммирует внутренние состояния в петле. Второй усредняет по сектору. Третий относится к одной нормированной частице. Голономный период \(\oint A\) является отдельной топологической величиной и вообще не является Hilbert trace.

Выбор между этими формами определяется типом observable, а не требуемым числом.

\section{Family example}

В четырёхмерном spin-menu projectors
\[
P_1=\frac14J,
\qquad
P_3=I-\frac14J
\]
имеют ранги один и три. Поэтому loop trace видит
\[
\operatorname{Tr}P_1=1,
\qquad
\operatorname{Tr}P_3=3,
\]
а normalized sector average даёт единичный вес в обоих случаях.

\begin{proposition}[Различие multiplicity и state normalization]
Один и тот же parent trace без свободного коэффициента даёт вес три для family-loop и вес один для нормированного family-state. Различие определяется observable type.
\end{proposition}

Это объясняет, почему требования ``везде использовать одну норму'' и ``везде сохранять rank'' одинаково неверны.

\section{Gauge example из \(SU(5)\)}

В fundamental representation
\[
Y=\operatorname{diag}\left(-\frac13,-\frac13,-\frac13,
\frac12,\frac12\right)
\]
имеет
\[
\operatorname{Tr}_{\mathbf5}Y^2=\frac56.
\]
Для стандартно нормированного nonabelian generator
\[
\operatorname{Tr}_{\mathbf5}T^2=\frac12.
\]
Поэтому parent \(SU(5)\)-trace фиксирует
\[
k_Y=\frac{\operatorname{Tr}Y^2}{\operatorname{Tr}T^2}
=\frac53
\]
и, при равенстве unified couplings,
\[
\sin^2\theta_W(\Lambda)=\frac38.
\]

Для полного chiral package \(\mathbf{10}\oplus\overline{\mathbf5}\) subgroup indices равны
\[
I_{SU(3)}=I_{SU(2)}=I_{U(1)}^{GUT}=2.
\]
Здесь разные subgroup traces действительно являются разными тенями одного representation trace, а их относительные веса не подгоняются.

\section{Почему это не спасает старую \(\tau\)-формулу}

Ненормированная constant mode имеет нормы
\[
\|1_{\mathbb{RP}^3}\|^2=\pi^2,
\qquad
\|1_{S^1}\|^2=2\pi.
\]
Но для нормированных particle states
\[
\int_{\mathbb{RP}^3}|\widehat1|^2=1,
\qquad
\int_{S^1}|\widehat1|^2=1.
\]
Следовательно, raw volumes допустимы в background action или loop sum, но не возникают автоматически в mass matrix element одной частицы. Вторая гипотеза объясняет прежний конфликт категорий, но не возвращает seed \(\pi^2+2\pi+2/3\).

\section{Вердикт}

Вторая сумасбродная гипотеза проходит в ограниченной форме:
\begin{quote}
Разные эффективные меры допустимы, если они получены из одного parent trace фиксированными projectors и заранее определённым типом observable.
\end{quote}

Она закрывается отрицательно в форме ``каждому сектору свой trace''. Она положительно объясняет family multiplicity и \(SU(5)\) gauge normalization, но пока не даёт нового low-energy prediction.

Следующий gate должен собрать menu projector, rank-one family operator и \(SU(5)\) representation trace в одной tensor-product algebra и проверить, действительно ли одна normalized parent state порождает обе структуры без смены соглашений.

\end{document}